Una lupa, unos lentes de contacto o una gota de agua sobre una hoja tienen algo en común: todas son superficies curvas que doblan la luz.

En la práctica de la ley de Snell la luz cruzaba una frontera plana. Ahí todos los rayos que llegaban con el mismo ángulo se doblaban exactamente igual, y un grupo de rayos paralelos seguía siendo paralelo después de entrar al acrílico. Pero si la frontera es curva la historia cambia por completo, porque la normal ya no apunta en la misma dirección en todos los puntos de la superficie, y entonces cada rayo llega con un ángulo distinto, se dobla una cantidad distinta y, si la curva tiene la forma adecuada, todos terminan juntándose en una misma región.

Esa es, en el fondo, la idea de una lente.

En esta práctica hicimos pasar varios rayos paralelos a través de un cilindro de acrílico y observamos cómo se acercaban entre sí, dónde se cruzaban y qué tan bien lo hacían. No tomamos medidas; el objetivo fue mirar con cuidado y explicar lo que veíamos con la teoría.
